\documentclass[aspectratio=169,10pt]{beamer}

% Shared Beamer setup for Fall 2026 course slides.
% Clean Cal Poly-inspired theme: no custom class files or uncommon packages.
\mode<presentation>{
  \usetheme{default}
  \usefonttheme{professionalfonts}
}

\definecolor{CalPolyGreen}{HTML}{154734}
\definecolor{CalPolyGold}{HTML}{BD8B13}
\definecolor{CalPolySage}{HTML}{7FA28F}
\definecolor{CalPolyMint}{HTML}{DDF1E7}
\definecolor{CalPolySand}{HTML}{A29F86}
\definecolor{CalPolyBeige}{HTML}{EFEEDE}
\definecolor{DeckInk}{HTML}{1F2933}
\definecolor{DeckMuted}{HTML}{5F6B7A}
\definecolor{DeckSoft}{HTML}{F7F7F5}

\setbeamersize{text margin left=0.55in,text margin right=0.55in}
\setbeamercolor{normal text}{fg=DeckInk,bg=white}
\setbeamercolor{structure}{fg=CalPolyGreen}
\setbeamercolor{title}{fg=white}
\setbeamercolor{frametitle}{fg=CalPolyGreen,bg=white}
\setbeamercolor{block title}{bg=CalPolySage,fg=white}
\setbeamercolor{block body}{bg=CalPolyMint,fg=black}
\setbeamercolor{alerted text}{fg=CalPolyGold}

\setbeamerfont{title}{size=\Huge,series=\bfseries}
\setbeamerfont{frametitle}{size=\Large,series=\bfseries}
\setbeamerfont{block title}{size=\normalsize,series=\bfseries}

\setbeamertemplate{navigation symbols}{}
\setbeamertemplate{itemize item}{\color{CalPolyGold}$\blacktriangleright$}
\setbeamertemplate{itemize subitem}{\color{CalPolyGreen}$\bullet$}
\setbeamertemplate{footline}{
  \hfill{\color{DeckMuted}\scriptsize\insertframenumber/\inserttotalframenumber}
  \hspace{0.18in}\vspace{0.12in}
}
\setbeamertemplate{frametitle}{
  \vspace{0.18in}
  \hspace{0.02in}\insertframetitle\par
  \vspace{0.08in}
  \noindent\makebox[\linewidth][l]{%
    {\color{CalPolyGold}\rule{0.55in}{2pt}}%
    {\color{CalPolyGreen}\rule{\dimexpr\linewidth-0.55in\relax}{2pt}}%
  }
}

\newcommand{\CourseNumber}{Course}
\newcommand{\CourseTitle}{Course Title}
\newcommand{\LectureNumber}{0}
\newcommand{\LectureTitle}{Lecture Title}
\newcommand{\LectureDate}{Date}
\newcommand{\CourseUnit}{Unit}

\newcommand{\coursenumber}[1]{\renewcommand{\CourseNumber}{#1}}
\newcommand{\coursetitle}[1]{\renewcommand{\CourseTitle}{#1}}
\newcommand{\lecturenumber}[1]{\renewcommand{\LectureNumber}{#1}}
\newcommand{\lecturetitle}[1]{\renewcommand{\LectureTitle}{#1}}
\newcommand{\lecturedate}[1]{\renewcommand{\LectureDate}{#1}}
\newcommand{\courseunit}[1]{\renewcommand{\CourseUnit}{#1}}

\author{Kirk Duran}
\institute{California Polytechnic State University, San Luis Obispo}
\date{\LectureDate}

\newcommand{\makedecktitle}{
  \begingroup
  \setbeamercolor{background canvas}{bg=CalPolyGreen}
  \begin{frame}[plain]
    \vfill
    \begin{center}
      {\usebeamerfont{title}\color{white}\LectureTitle\par}
      \vspace{0.18in}
      {\Large\color{white}Lecture \LectureNumber\par}
    \end{center}
    \vfill
    \noindent{\color{white}\rule{\linewidth}{1.2pt}}\par
    \vspace{0.08in}
    \noindent\colorbox{CalPolyGold}{\makebox[0.24\linewidth][c]{\color{white}\Large\CourseNumber}}
    \hspace{0.12in}{\color{white}\Large Kirk Duran}
  \end{frame}
  \endgroup
}

\newenvironment{keyidea}{\begin{block}{Key idea}}{\end{block}}
\newenvironment{activity}{\begin{block}{In class}}{\end{block}}
\newenvironment{cpdefinition}{\begin{block}{Definition}}{\end{block}}
\newenvironment{exampleblockcp}{
  \begingroup
  \setbeamercolor{block title}{bg=CalPolySand,fg=white}
  \setbeamercolor{block body}{bg=CalPolyBeige,fg=black}
  \begin{block}{Example}
}{
  \end{block}
  \endgroup
}

\newcommand{\imageplaceholder}[2][0.3in]{%
  \noindent\makebox[\linewidth][c]{%
    \fcolorbox{CalPolySage}{CalPolyBeige}{%
      \parbox[c][#1][c]{0.86\linewidth}{%
        \centering
        {\color{DeckMuted}\scriptsize Image placeholder: }%
        {\color{DeckInk}\scriptsize\ttfamily\detokenize{#2}}%
      }%
    }%
  }%
}

\newcommand{\sectionframe}[1]{
  \begingroup
  \setbeamercolor{background canvas}{bg=CalPolyGreen}
  \begin{frame}[plain]
    \vfill
    {\color{white}\LARGE\bfseries #1\par}
    \vspace{0.18in}
    {\color{CalPolyGold}\rule{0.22\paperwidth}{3pt}}
    {\color{white}\rule{0.62\paperwidth}{3pt}}
    \vfill
  \end{frame}
  \endgroup
}

\newcommand{\placeholdernote}{
  \begin{block}{Placeholder}
    Replace these bullets with the final lecture material, examples, and in-class activities.
  \end{block}
}


\usepackage{tikz}

\graphicspath{{../assets/lecture-13/}}

% If an image file is not yet available, skip the visual slot.
\newcommand{\lectureimage}[2][1.75in]{%
  \IfFileExists{../assets/lecture-13/#2}{%
    \includegraphics[width=\linewidth,height=#1,keepaspectratio]{#2}%
  }{}%
}

% Explicit placeholder for visuals/examples that will be added later.
\newcommand{\lectureplaceholder}[2][2.35in]{%
  \begin{center}
    \fbox{%
      \begin{minipage}[c][#1][c]{0.90\linewidth}
        \centering
        \small\textit{#2}
      \end{minipage}%
    }
  \end{center}%
}

\setbeamersize{text margin left=0.42in,text margin right=0.42in}

\coursenumber{CS 4880}
\coursetitle{Artificial Intelligence}
\lecturenumber{13}
\lecturetitle{Propositional Logic: Logical Agents, Syntax, Semantics, and Entailment}
\lecturedate{September 23, 2026}
\courseunit{Logic and Knowledge Representation}

\begin{document}

% Estimated pacing: 50 minutes
% Motivation + logical agents: 10 min
% Historical context: 7 min
% Propositional syntax: 12 min
% Semantics + models: 11 min
% Entailment: 8 min
% Wrap / bridge: 2 min

% -----------------------------------------------------------------------------
% Slide 1
% -----------------------------------------------------------------------------
\makedecktitle

\section{Why Logic in AI?}

% -----------------------------------------------------------------------------
% Slide 2
% -----------------------------------------------------------------------------
\begin{frame}[shrink=10]{A Different Kind of AI Question}
  \begin{columns}[T,onlytextwidth]
    \begin{column}{0.53\textwidth}
      Imagine a warehouse robot approaching a charging bay.

      \begin{block}{What the robot knows}
        \[
          \mathrm{BatteryLow}
        \]
        \[
          \mathrm{SmokeDetected} \rightarrow \mathrm{BayUnsafe}
        \]
        \[
          \mathrm{BatteryLow}\land\neg\mathrm{BayUnsafe}
          \rightarrow \mathrm{EnterBay}
        \]
      \end{block}

      The immediate question is not ``Which path is shortest?''

      \begin{center}
        \textbf{What conclusions are justified by what I know?}
      \end{center}
    \end{column}
    \begin{column}{0.41\textwidth}
      \lectureimage[2.75in]{slide2.png}
    \end{column}
  \end{columns}

  \begin{keyidea}
    Search reasons over possible actions. Logic reasons over what must be true given an agent's knowledge.
  \end{keyidea}
\end{frame}

% -----------------------------------------------------------------------------
% Slide 3
% -----------------------------------------------------------------------------
\begin{frame}[shrink=10]{From Search Agents to Logical Agents}
  \begin{columns}[T,onlytextwidth]
    \begin{column}{0.47\textwidth}
      \begin{block}{Search-oriented agent}
        \begin{enumerate}
          \item Represent a state.
          \item Generate successor states.
          \item Evaluate paths or actions.
          \item Choose an action.
        \end{enumerate}
      \end{block}
    \end{column}
    \begin{column}{0.47\textwidth}
      \begin{block}{Logical agent}
        \begin{enumerate}
          \item Represent facts and rules.
          \item Add observations to a knowledge base.
          \item Infer consequences.
          \item Act using those conclusions.
        \end{enumerate}
      \end{block}
    \end{column}
  \end{columns}

  \lectureimage[1.35in]{slide3.png}
\end{frame}

% -----------------------------------------------------------------------------
% Slide 4
% -----------------------------------------------------------------------------
\begin{frame}[shrink=10]{The Knowledge-Based Agent Loop}
  \begin{columns}[T,onlytextwidth]
    \begin{column}{0.54\textwidth}
      A classical knowledge-based agent repeatedly updates and queries an explicit knowledge base.

      \begin{enumerate}
        \item \textbf{Perceive} something about the world.
        \item \textbf{TELL} the knowledge base what was observed.
        \item \textbf{ASK} what follows from the current knowledge.
        \item Choose an action.
        \item Record the action and repeat.
      \end{enumerate}

      \begin{block}{Two separate problems}
        \textbf{Representation:} What can the agent say?\\[0.04in]
        \textbf{Inference:} What can the agent conclude?
      \end{block}
    \end{column}
    \begin{column}{0.40\textwidth}
      \lectureimage[2.75in]{slide4.png}
    \end{column}
  \end{columns}
\end{frame}

% -----------------------------------------------------------------------------
% Slide 5
% -----------------------------------------------------------------------------
\begin{frame}[shrink=10]{A Knowledge Base Stores More Than Raw Data}
  \begin{columns}[T,onlytextwidth]
    \begin{column}{0.53\textwidth}
      A database might store isolated facts:

      \[
        \mathrm{BatteryLow}
      \]
      \[
        \mathrm{ChargerAvailable}
      \]

      A knowledge base may also store \textbf{rules connecting facts}:

      \[
        \mathrm{BatteryLow}\land\mathrm{ChargerAvailable}
        \rightarrow \mathrm{ShouldCharge}
      \]

      The sentence \(\mathrm{ShouldCharge}\) does not need to be written explicitly if it can be derived.
    \end{column}
    \begin{column}{0.41\textwidth}
      \lectureimage[2.55in]{slide5.png}
    \end{column}
  \end{columns}

  \begin{keyidea}
    Inference turns explicit knowledge into implicit knowledge.
  \end{keyidea}
\end{frame}

% -----------------------------------------------------------------------------
% Slide 6
% -----------------------------------------------------------------------------
\begin{frame}[shrink=10]{Explicit Knowledge vs. Implicit Knowledge}
  \begin{columns}[T,onlytextwidth]
    \begin{column}{0.47\textwidth}
      \begin{block}{Explicitly stored}
        \[
          P
        \]
        \[
          P\rightarrow Q
        \]
      \end{block}

      These sentences are directly present in the knowledge base.
    \end{column}
    \begin{column}{0.47\textwidth}
      \begin{block}{Implicitly supported}
        \[
          Q
        \]
      \end{block}

      The agent may be justified in accepting \(Q\), even if \(Q\) was never entered directly.
    \end{column}
  \end{columns}

  \begin{keyidea}
    A logical agent is valuable because it can derive consequences that were not explicitly enumerated in advance.
  \end{keyidea}
\end{frame}

% -----------------------------------------------------------------------------
% Slide 7
% -----------------------------------------------------------------------------
\begin{frame}[shrink=10]{Unknown Is Not the Same as False}
  \begin{columns}[T,onlytextwidth]
    \begin{column}{0.53\textwidth}
      Suppose the knowledge base contains
      \[
        \mathrm{SmokeDetected}\rightarrow\mathrm{BayUnsafe}
      \]
      but contains no fact telling us whether smoke is currently detected.

      Can the robot conclude
      \[
        \neg\mathrm{BayUnsafe}?
      \]

      \begin{alertblock}{No.}
        ``The KB does not prove \(P\)'' is different from ``the KB proves \(\neg P\).''
      \end{alertblock}
    \end{column}
    \begin{column}{0.41\textwidth}
      \lectureimage[2.55in]{slide7.png}
    \end{column}
  \end{columns}
\end{frame}

\section{Historical Context: Symbolic AI}

% -----------------------------------------------------------------------------
% Slide 8
% -----------------------------------------------------------------------------
\begin{frame}[shrink=10]{Good Old-Fashioned AI: Make Knowledge Explicit}
  \begin{columns}[T,onlytextwidth]
    \begin{column}{0.55\textwidth}
      Much early AI research treated intelligence as the manipulation of explicit symbolic representations.

      \begin{itemize}
        \item Facts were represented as symbols.
        \item Rules specified how symbols could be transformed.
        \item Reasoning produced new symbolic statements.
        \item The internal reasoning process was often inspectable.
      \end{itemize}

      The later label \textbf{GOFAI}---``Good Old-Fashioned Artificial Intelligence''---is commonly used for this symbolic tradition.
    \end{column}
    \begin{column}{0.39\textwidth}
      \lectureimage[2.70in]{slide8.png}
    \end{column}
  \end{columns}

  \begin{keyidea}
    The central wager was that intelligent behavior could emerge from explicit knowledge plus general reasoning procedures.
  \end{keyidea}
\end{frame}

% -----------------------------------------------------------------------------
% Slide 9
% -----------------------------------------------------------------------------
\begin{frame}[shrink=10]{Logic Theorist: Reasoning as Computation}
  \begin{columns}[T,onlytextwidth]
    \begin{column}{0.53\textwidth}
      Newell, Simon, and Shaw's \textit{Logic Theorist} in the 1950s demonstrated that a program could search for formal proofs in symbolic logic.

      \begin{itemize}
        \item The program represented mathematical statements symbolically.
        \item It applied legal transformation rules.
        \item It searched through possible proof steps.
      \end{itemize}

      Notice the recurring AI pattern:
      \begin{center}
        \textbf{representation + operators + search}
      \end{center}
    \end{column}
    \begin{column}{0.41\textwidth}
      \lectureimage[2.65in]{slide9.png}
    \end{column}
  \end{columns}
\end{frame}

% -----------------------------------------------------------------------------
% Slide 10
% -----------------------------------------------------------------------------
\begin{frame}[shrink=10]{Expert Systems: Put the Knowledge in Rules}
  \begin{columns}[T,onlytextwidth]
    \begin{column}{0.53\textwidth}
      Later expert systems made the architecture especially explicit:

      \begin{block}{Knowledge base}
        Domain-specific facts and IF--THEN rules.
      \end{block}

      \begin{block}{Inference engine}
        A more general mechanism that decides which rules apply and derives conclusions.
      \end{block}

      This separation is still conceptually important today: \textbf{what the system knows} is different from \textbf{how it reasons}.
    \end{column}
    \begin{column}{0.41\textwidth}
      \lectureimage[2.65in]{slide10.png}
    \end{column}
  \end{columns}
\end{frame}

% -----------------------------------------------------------------------------
% Slide 11
% -----------------------------------------------------------------------------
\begin{frame}[shrink=10]{Why Symbolic AI Was Powerful---and Why It Was Hard}
  \begin{columns}[T,onlytextwidth]
    \begin{column}{0.47\textwidth}
      \begin{block}{Strengths}
        \begin{itemize}
          \item explicit, inspectable knowledge;
          \item precise rules of reasoning;
          \item strong performance in structured domains;
          \item conclusions can often be explained as proofs.
        \end{itemize}
      \end{block}
    \end{column}
    \begin{column}{0.47\textwidth}
      \begin{block}{Difficulties}
        \begin{itemize}
          \item knowledge engineering is expensive;
          \item the real world is messy and incomplete;
          \item rule sets can become large and brittle;
          \item inference itself can become computationally expensive.
        \end{itemize}
      \end{block}
    \end{column}
  \end{columns}

  \begin{keyidea}
    Symbolic AI did not disappear; logic remains one of the clearest ways to study representation, reasoning, and guarantees.
  \end{keyidea}
\end{frame}

\section{Propositional Logic: The Language}

% -----------------------------------------------------------------------------
% Slide 12
% -----------------------------------------------------------------------------
\begin{frame}[shrink=10]{Propositional Logic: The Simplest Useful Logic}
  \begin{cpdefinition}
    A \textbf{proposition} is a statement that is either true or false.
  \end{cpdefinition}

  \begin{columns}[T,onlytextwidth]
    \begin{column}{0.49\textwidth}
      Examples:
      \[
        R = \text{``It is raining.''}
      \]
      \[
        B = \text{``The battery is low.''}
      \]
      \[
        U = \text{``The charging bay is unsafe.''}
      \]
    \end{column}
    \begin{column}{0.45\textwidth}
      Propositional logic treats each atomic statement as a single indivisible symbol.

      It does \textbf{not} yet represent objects, properties, variables, or quantifiers.
    \end{column}
  \end{columns}

  \begin{keyidea}
    Propositional logic is limited, but its simplicity makes the mechanics of logical reasoning very clear.
  \end{keyidea}
\end{frame}

% -----------------------------------------------------------------------------
% Slide 13
% -----------------------------------------------------------------------------
\begin{frame}[shrink=10]{Atomic Symbols Need an Interpretation}
  \begin{columns}[T,onlytextwidth]
    \begin{column}{0.52\textwidth}
      The symbol \(P\) has no meaning by itself.

      We assign it an interpretation such as
      \[
        P = \text{``Package 17 has arrived.''}
      \]

      Another symbol might mean
      \[
        Q = \text{``Robot 2 is available.''}
      \]

      Logic manipulates the symbols, while the interpretation connects those symbols to a domain.
    \end{column}
    \begin{column}{0.42\textwidth}
      \lectureimage[2.50in]{slide13-interpretation.png}
    \end{column}
  \end{columns}
\end{frame}

% -----------------------------------------------------------------------------
% Slide 14
% -----------------------------------------------------------------------------
\begin{frame}[shrink=8]{Building Sentences with Logical Connectives}
  \begin{center}
    \renewcommand{\arraystretch}{1.35}
    \begin{tabular}{c|l|l}
      \textbf{Symbol} & \textbf{Name} & \textbf{Example} \\\hline
      $\neg$ & NOT / negation & $\neg P$ \\
      $\land$ & AND / conjunction & $P\land Q$ \\
      $\lor$ & OR / disjunction & $P\lor Q$ \\
      $\rightarrow$ & implication & $P\rightarrow Q$ \\
      $\leftrightarrow$ & biconditional & $P\leftrightarrow Q$ \\
    \end{tabular}
  \end{center}

  \begin{block}{Example}
    \[
      \mathrm{BatteryLow}\land\neg\mathrm{BayUnsafe}
      \rightarrow\mathrm{EnterBay}
    \]
  \end{block}

  Complex sentences are built recursively from atomic propositions and connectives.
\end{frame}

% -----------------------------------------------------------------------------
% Slide 15
% -----------------------------------------------------------------------------
\begin{frame}[shrink=10]{Parentheses Matter}
  Compare
  \[
    (P\lor Q)\land R
  \]
  with
  \[
    P\lor(Q\land R).
  \]

  These are different sentences and can have different truth values in the same world.

  \begin{block}{Practical rule}
    When in doubt, use parentheses. A logic expression should communicate its structure unambiguously to both the reader and the inference system.
  \end{block}

  \lectureimage[1.55in]{slide15.png}
\end{frame}

% -----------------------------------------------------------------------------
% Slide 16
% -----------------------------------------------------------------------------
\begin{frame}[shrink=10]{A Knowledge Base Is a Set of Sentences}
  Suppose a delivery robot's KB contains

  \[
    B
  \]
  \[
    B\rightarrow C
  \]
  \[
    C\rightarrow M
  \]

  where
  \begin{itemize}
    \item \(B\): battery is low;
    \item \(C\): robot should charge;
    \item \(M\): robot should move toward the charger.
  \end{itemize}

  \begin{keyidea}
    The knowledge base is not one fact; it is a collection of sentences whose consequences interact.
  \end{keyidea}
\end{frame}

% -----------------------------------------------------------------------------
% Slide 17
% -----------------------------------------------------------------------------
\begin{frame}[shrink=10]{What Does $P\rightarrow Q$ Actually Say?}
  \begin{columns}[T,onlytextwidth]
    \begin{column}{0.54\textwidth}
      An implication
      \[
        P\rightarrow Q
      \]
      states that whenever \(P\) is true, \(Q\) must also be true.

      The implication is violated only by a situation where
      \[
        P=\mathrm{true},\qquad Q=\mathrm{false}.
      \]

      Example:
      \[
        \mathrm{BatteryLow}\rightarrow\mathrm{ShouldCharge}.
      \]
    \end{column}
    \begin{column}{0.40\textwidth}
      \lectureimage[2.50in]{slide17.png}
    \end{column}
  \end{columns}
\end{frame}

% -----------------------------------------------------------------------------
% Slide 18
% -----------------------------------------------------------------------------
\begin{frame}[shrink=10]{Implication Is Not Causation}
  \begin{columns}[T,onlytextwidth]
    \begin{column}{0.52\textwidth}
      The sentence
      \[
        \mathrm{Rain}\rightarrow\mathrm{WetGround}
      \]
      does \textbf{not} claim that rain is the only possible cause of wet ground.

      Sprinklers might also make the ground wet.

      Propositional implication expresses a truth constraint:
      \begin{center}
        \textit{if the antecedent holds, the consequent must hold.}
      \end{center}
    \end{column}
    \begin{column}{0.42\textwidth}
      \lectureimage[2.55in]{slide18.png}
    \end{column}
  \end{columns}

  \begin{keyidea}
    Logical implication describes a relationship among truth values, not a physical mechanism.
  \end{keyidea}
\end{frame}

% -----------------------------------------------------------------------------
% Slide 19
% -----------------------------------------------------------------------------
\begin{frame}[shrink=8]{Do Not Run the Rule Backward}
  Suppose the KB contains
  \[
    \mathrm{Rain}\rightarrow\mathrm{WetGround}
  \]
  and we observe
  \[
    \mathrm{WetGround}.
  \]

  Can we conclude \(\mathrm{Rain}\)?

  \begin{alertblock}{No.}
    The implication only guarantees wet ground when rain occurs. It does not say wet ground can only be caused by rain.
  \end{alertblock}

  \begin{block}{Likewise}
    From \(P\rightarrow Q\) and \(Q\), we cannot generally conclude \(P\).
  \end{block}
\end{frame}

% -----------------------------------------------------------------------------
% Slide 20
% -----------------------------------------------------------------------------
\begin{frame}[shrink=8]{Quick Check: Which Sentences Say What We Mean?}
  Let
  \[
    S=\text{``smoke is detected''},\qquad U=\text{``the bay is unsafe.''}
  \]

  Which sentence best represents
  \begin{center}
    ``If smoke is detected, then the bay is unsafe''?
  \end{center}

  \begin{enumerate}
    \item $S\land U$
    \item $S\rightarrow U$
    \item $U\rightarrow S$
    \item $S\leftrightarrow U$
  \end{enumerate}

  \begin{block}{Answer}
    \(S\rightarrow U\). The other formulas make stronger or simply different claims.
  \end{block}
\end{frame}

\section{Semantics: What Do the Sentences Mean?}

% -----------------------------------------------------------------------------
% Slide 21
% -----------------------------------------------------------------------------
\begin{frame}[shrink=10]{Syntax vs. Semantics}
  \begin{columns}[T,onlytextwidth]
    \begin{column}{0.47\textwidth}
      \begin{block}{Syntax}
        The legal structure of expressions.

        \[
          (P\land Q)\rightarrow R
        \]

        Syntax answers:
        \begin{center}
          \textit{Is this a well-formed sentence?}
        \end{center}
      \end{block}
    \end{column}
    \begin{column}{0.47\textwidth}
      \begin{block}{Semantics}
        The conditions under which a sentence is true or false.

        Semantics answers:
        \begin{center}
          \textit{What would the world have to be like for this sentence to be true?}
        \end{center}
      \end{block}
    \end{column}
  \end{columns}

  \begin{keyidea}
    Logic needs both a formal language and a precise account of what sentences mean.
  \end{keyidea}
\end{frame}

% -----------------------------------------------------------------------------
% Slide 22
% -----------------------------------------------------------------------------
\begin{frame}[shrink=10]{A Model Is One Possible World}
  \begin{cpdefinition}
    A \textbf{model} assigns a truth value to every proposition symbol under consideration.
  \end{cpdefinition}

  With two symbols \(P\) and \(Q\), one model might be
  \[
    P=\mathrm{true},\qquad Q=\mathrm{false}.
  \]

  Another model might be
  \[
    P=\mathrm{false},\qquad Q=\mathrm{false}.
  \]

  Each model represents one logically possible configuration of the world at the level of detail captured by those symbols.
\end{frame}

% -----------------------------------------------------------------------------
% Slide 23
% -----------------------------------------------------------------------------
\begin{frame}[shrink=10]{Models Grow Exponentially}
  With \(n\) independent proposition symbols, there are
  \[
    2^n
  \]
  possible truth assignments.

  \begin{center}
    \renewcommand{\arraystretch}{1.25}
    \begin{tabular}{c|c}
      \textbf{Symbols} & \textbf{Possible models} \\\hline
      1 & 2 \\
      2 & 4 \\
      3 & 8 \\
      10 & 1,024 \\
    \end{tabular}
  \end{center}

  \begin{keyidea}
    Thinking in models gives logic a precise semantics. Next lecture, the number of models becomes a computational problem.
  \end{keyidea}
\end{frame}

% -----------------------------------------------------------------------------
% Slide 24
% -----------------------------------------------------------------------------
\begin{frame}[shrink=10]{Satisfaction: Does a Model Make a Sentence True?}
  Suppose model \(m\) assigns
  \[
    P=\mathrm{true},\qquad Q=\mathrm{false}.
  \]

  Then
  \[
    m\models P
  \]
  and
  \[
    m\models \neg Q.
  \]

  But
  \[
    m\not\models P\land Q.
  \]

  \begin{cpdefinition}
    We write \(m\models\alpha\) when model \(m\) \textbf{satisfies} sentence \(\alpha\).
  \end{cpdefinition}
\end{frame}

% -----------------------------------------------------------------------------
% Slide 25
% -----------------------------------------------------------------------------
\begin{frame}[shrink=10]{A Knowledge Base Rules Out Possible Worlds}
  \begin{columns}[T,onlytextwidth]
    \begin{column}{0.53\textwidth}
      Suppose
      \[
        KB=\{P,\;P\rightarrow Q\}.
      \]

      Any model in which \(P\) is false is incompatible with the first sentence.

      Any model in which \(P\) is true and \(Q\) is false is incompatible with the second sentence.

      The models that survive are exactly the worlds consistent with \textbf{all} sentences in the KB.
    \end{column}
    \begin{column}{0.41\textwidth}
      \lectureimage[2.70in]{slide25.png}
    \end{column}
  \end{columns}

  \begin{keyidea}
    Knowledge constrains the set of worlds the agent considers possible.
  \end{keyidea}
\end{frame}

\section{Entailment: What Must Follow?}

% -----------------------------------------------------------------------------
% Slide 26
% -----------------------------------------------------------------------------
\begin{frame}[shrink=8]{Entailment}
  \begin{cpdefinition}
    \[
      KB\models\alpha
    \]
    means that \(\alpha\) is true in \textbf{every model} in which the knowledge base is true.
  \end{cpdefinition}

  Read it as:
  \begin{center}
    ``The knowledge base \textbf{entails} \(\alpha\).''
  \end{center}

  Or more intuitively:
  \begin{center}
    \textbf{There is no world consistent with what I know in which \(\alpha\) is false.}
  \end{center}

  \begin{keyidea}
    Entailment is about necessity relative to the knowledge base, not about whether a conclusion merely seems plausible.
  \end{keyidea}
\end{frame}

% -----------------------------------------------------------------------------
% Slide 27
% -----------------------------------------------------------------------------
\begin{frame}[shrink=10]{Entailment as Model-Set Containment}
  Let \(M(KB)\) be the set of models satisfying the knowledge base and \(M(\alpha)\) the set satisfying \(\alpha\).

  Then
  \[
    KB\models\alpha
  \]
  exactly when
  \[
    M(KB)\subseteq M(\alpha).
  \]

  \lectureimage[2.10in]{slide27.png}
\end{frame}

% -----------------------------------------------------------------------------
% Slide 28
% -----------------------------------------------------------------------------
\begin{frame}[shrink=10]{Example: Does the KB Entail the Conclusion?}
  Consider
  \[
    KB=\{\mathrm{BatteryLow},\;
    \mathrm{BatteryLow}\rightarrow\mathrm{ShouldCharge}\}.
  \]

  Query:
  \[
    \mathrm{ShouldCharge}?
  \]

  Ask the semantic question:
  \begin{center}
    Is there \textbf{any} model in which both KB sentences are true but \(\mathrm{ShouldCharge}\) is false?
  \end{center}

  There is not. Therefore
  \[
    KB\models\mathrm{ShouldCharge}.
  \]

  \begin{block}{Important}
    Today we are describing \textbf{why} the entailment holds. Next lecture we turn that idea into an inference procedure.
  \end{block}
\end{frame}

% -----------------------------------------------------------------------------
% Slide 29
% -----------------------------------------------------------------------------
\begin{frame}[shrink=10]{Non-Entailment Needs Only One Counterexample}
  Suppose
  \[
    KB=\{\mathrm{WetGround}\}.
  \]

  Does
  \[
    KB\models\mathrm{Rain}
  \]
  hold?

  No. A model with
  \[
    \mathrm{WetGround}=\mathrm{true},\qquad
    \mathrm{Rain}=\mathrm{false}
  \]
  is perfectly consistent with the KB.

  \begin{keyidea}
    To show entailment, every KB-model must support the conclusion. To refute entailment, one countermodel is enough.
  \end{keyidea}
\end{frame}

% -----------------------------------------------------------------------------
% Slide 30
% -----------------------------------------------------------------------------
\begin{frame}[shrink=10]{Classic Logical-Agent Example: Wumpus World}
  \begin{columns}[T,onlytextwidth]
    \begin{column}{0.54\textwidth}
      In the classic Wumpus World, an agent receives local percepts such as breezes and stenches but must reason about hidden hazards.

      A rule might encode an idea like
      \[
        B_{1,1}\rightarrow(P_{1,2}\lor P_{2,1}),
      \]
      meaning a breeze at one square implies a pit in at least one neighboring square.

      The agent's job is to combine percepts and rules to determine which locations are safe or dangerous.
    \end{column}
    \begin{column}{0.40\textwidth}
      \lectureimage[2.70in]{slide30.png}
    \end{column}
  \end{columns}

  \begin{keyidea}
    A logical agent can reason about states of the world it has never directly observed.
  \end{keyidea}
\end{frame}

% -----------------------------------------------------------------------------
% Slide 31
% -----------------------------------------------------------------------------
\begin{frame}[shrink=8]{What We Have---and What We Still Need}
  We now have the pieces needed to state a logical reasoning problem:

  \begin{enumerate}
    \item A formal language for writing sentences.
    \item A knowledge base containing sentences.
    \item Models that define what those sentences mean.
    \item Entailment to define what conclusions logically follow.
  \end{enumerate}

  But there is still a practical question:

  \begin{center}
    \Large\textbf{How does the computer determine whether $KB\models\alpha$?}
  \end{center}

  \begin{block}{Next lecture: Propositional Inference}
    Logical equivalence, truth tables, model checking, and inference procedures.
  \end{block}
\end{frame}

% -----------------------------------------------------------------------------
% Slide 32
% -----------------------------------------------------------------------------
\begin{frame}[shrink=8]{Takeaways}
  \begin{itemize}
    \item Knowledge-based agents explicitly represent facts and rules, then derive consequences through inference.
    \item Propositional logic builds complex sentences from atomic true/false propositions and logical connectives.
    \item \textbf{Syntax} determines legal sentence structure; \textbf{semantics} determines truth in a model.
    \item A knowledge base constrains which possible worlds remain consistent with what the agent knows.
    \item \(KB\models\alpha\) means \(\alpha\) is true in every model satisfying the knowledge base.
    \item Entailment defines what logically follows; an inference algorithm tells us how to compute or prove it.
  \end{itemize}

  \begin{keyidea}
    Logic gives an AI agent an explicit language for knowledge and a precise standard for when a conclusion is justified.
  \end{keyidea}
\end{frame}

\end{document}
